% SPDX-License-Identifier: MIT
% Copyright (c) 2026 handsomeZR-netizen
%
% Original CUMCM 2026 electronic-paper assembly template for mathmodel-skill.
% This is a rules-aligned repository template, not an official contest template.
% It intentionally omits the commitment form, numbering page, and identity data.

\documentclass[UTF8,12pt,a4paper,fontset=fandol]{ctexart}

\usepackage[margin=2.5cm]{geometry}
\usepackage{amsmath,amssymb,amsthm}
\usepackage{graphicx}
\usepackage{float}
\usepackage{booktabs}
\usepackage{longtable}
\usepackage{multirow}
\usepackage{listings}
\usepackage{xcolor}
\usepackage{hyperref}
\usepackage{url}
\usepackage{cite}
\usepackage{caption}
\usepackage{subcaption}
\usepackage{enumitem}

% pandoc 对**没有标题的表格**会生成 `{\def\LTcaptype{none} \begin{longtable}...}`，
% 本意是"别增加 table 计数器"。裸 longtable 下这没问题（不带 \caption 就不会用到
% \LTcaptype），但 caption 宏包会去 \refstepcounter{\LTcaptype}，于是报
% `! LaTeX Error: No counter 'none' defined.`，**而且 xelatex 照样输出一份残缺 PDF**
% （2026-09-03 实测：整篇论文只剩 1 页）。定义这个一次性计数器把它接住；
% 它与 table 计数器无关，真正的表号不会被打乱。
\newcounter{none}

\hypersetup{
  unicode=true,
  hidelinks,
  pdfauthor={},
  pdftitle={},
  pdfsubject={},
  pdfkeywords={}
}
\urlstyle{same}

\setlength{\parindent}{2em}
\setlength{\parskip}{0pt}
\linespread{1.25}
\captionsetup{font=small,labelfont=bf}
\lstset{
  basicstyle=\small\ttfamily,
  breaklines=true,
  frame=single,
  numbers=left,
  numberstyle=\tiny,
  showstringspaces=false,
  columns=fullflexible,
  keepspaces=true
}

% Pandoc may emit \tightlist in a LaTeX fragment.
\providecommand{\tightlist}{%
  \setlength{\itemsep}{0pt}\setlength{\parskip}{0pt}}

% These counters enforce the electronic-paper structure during compilation.
% The abstract occupies page 1. The body begins on page 2.
% 官方《论文格式规范》第四条的原文是"正文……**尽量控制在 20 页以内**"，
% 不是硬上限（上游把它写成 "limited to 30 pages" 是误读，见
% competitions/cumcm/current_rules.md 第 8 行）。所以这里 >20 页只**警告**，
% 留一个 >30 页的硬报错作远端兜底。参考文献与 AI 使用声明计入正文，附录不计。
\newcounter{mathmodelbodystartpage}
\newcounter{mathmodelbodypages}

\begin{document}

% MATHMODEL_* fields are fail-closed rendering tokens. render_paper.py fills
% them from decision_log.paper_metadata or explicit CLI arguments.
\pagenumbering{arabic}
\pagestyle{plain}
\thispagestyle{plain}

\begin{center}
  {\zihao{3}\bfseries MATHMODEL_CUMCM_TITLE}\par
  \vspace{1em}
  {\zihao{4}\bfseries 摘要}
\end{center}

% MATHMODEL:SECTION abstract

\par\noindent\textbf{关键词：}MATHMODEL_CUMCM_KEYWORDS

% The electronic paper must start with a one-page abstract sheet. Starting the
% body on any page other than page 2 is a submission-format error.
\clearpage
\ifnum\value{page}=2\relax\else
  \PackageError{mathmodel-cumcm}
    {CUMCM abstract must fit on the first page}
    {Shorten the title, abstract, or keywords before submission.}
\fi
\setcounter{mathmodelbodystartpage}{\value{page}}

% The Markdown section files own their headings. No table of contents is added.
% MATHMODEL:SECTION 1_problem_restate

% MATHMODEL:SECTION 2_problem_analysis

% MATHMODEL:SECTION 3_assumptions

% MATHMODEL:SECTION 4_notation

% MATHMODEL:SECTION 5_models

% MATHMODEL:SECTION 6_sensitivity

% MATHMODEL:SECTION 7_evaluation

% ---------------------------------------------------------------------------
% AI 工具使用声明 —— 《人工智能工具使用规定（2026 年试行）》要求它出现在
% **参考文献之前**，且两句话二选一、原文照抄。缺失或改写属于"故意隐瞒 / 虚假声明"，
% 取消评奖资格。render_ai_usage.py 按台账只生成其中一个源文件，
% 渲染器会把另一个 OPTIONAL 块整段删掉。
% （曾经放在参考文献之后，且只有"未使用"一种情况——那样用了 AI 的队伍
%   交出去的论文里一条声明都没有。已修正。）
% ---------------------------------------------------------------------------
\par\bigskip
% MATHMODEL:OPTIONAL cumcm_no_ai_statement BEGIN
\section*{AI 工具使用声明}
% MATHMODEL:SECTION cumcm_no_ai_statement
% MATHMODEL:OPTIONAL cumcm_no_ai_statement END
% MATHMODEL:OPTIONAL cumcm_ai_statement BEGIN
\section*{AI 工具使用声明}
% MATHMODEL:SECTION cumcm_ai_statement
% MATHMODEL:OPTIONAL cumcm_ai_statement END
\par\bigskip

% MATHMODEL:SECTION 8_references

% Flush deferred body floats before measuring. At this point \value{page} is
% the first appendix page, so subtracting the body start yields the body length.
\clearpage
\setcounter{mathmodelbodypages}{%
  \numexpr\value{page}-\value{mathmodelbodystartpage}\relax}
\typeout{MATHMODEL:CUMCM_BODY_PAGES=\arabic{mathmodelbodypages}}
\ifnum\value{mathmodelbodypages}>20\relax
  \PackageWarning{mathmodel-cumcm}
    {CUMCM body is \arabic{mathmodelbodypages} pages; the official format rule
     asks to keep the main text within 20 pages (appendices are counted
     separately). Consider moving derivations or long tables to the appendix}
\fi
\ifnum\value{mathmodelbodypages}>30\relax
  \PackageError{mathmodel-cumcm}
    {CUMCM body is far beyond the official "within 20 pages" guidance}
    {Shorten the main text. Appendices are counted separately.}
\fi

\appendix
% MATHMODEL:SECTION appendix_code

\end{document}
